% %%%%%%%%%%%%%%%%%%%%%%%%%%%%%%%%%%%%%%%%%%%%%%%%%%%%%%%%%%%
% \begin{frame}[fragile]\frametitle{}
% \begin{center}
% \Large{MidcurveNN : Encoder-Decoder Neural Network for Computing Midcurve of a Thin Polygon}
% \end{center}
% \end{frame}

%%%%%%%%%%%%%%%%%%%%%%%%%%%%%%%%%%%%%%%%%%%%%%%%%%%%%%%%%%%%%%%%%%%%%%%%%%%%%%%%%%
\begin{frame}[fragile]\frametitle{}
\begin{center}
{\Large Introduction To Midcurve}
\end{center}
\end{frame}

%%%%%%%%%%%%%%%%%%%%%%%%%%%%%%%%%%%%%%%%%%%%%%%%%%%%%%%%%%%%%%%%%%%%%%%%%%%%%%%%%%
\begin{frame}[fragile]\frametitle{}
\begin{center}
\includegraphics[width=0.9\linewidth,keepaspectratio]{midcurve1}
\end{center}
\end{frame}


%%%%%%%%%%%%%%%%%%%%%%%%%%%%%%%%%%%%%%%%%%%%%%%%%%%%%%%%%%%
\begin{frame}[fragile]\frametitle{Can we use shapes directly?}
	\begin{itemize}
	\item CAD : Designing Shapes
	\item CAE : Engineering Analysis
	\item CAD$\rightarrow$CAE: Simplification for quicker results.
	\end{itemize}

\begin{center}
\includegraphics[width=\linewidth,keepaspectratio]{midcurve2}
\end{center}

\end{frame}

%%%%%%%%%%%%%%%%%%%%%%%%%%%%%%%%%%%%%%%%%%%%%%%%%%%%%%%%%%%%%%%%%%%%%%%%%%%%%%%%%%
\begin{frame}[fragile]\frametitle{CAD-CAE}
\begin{center}
\includegraphics[width=0.9\linewidth,keepaspectratio]{midcurve3}
\end{center}
\end{frame}

%%%%%%%%%%%%%%%%%%%%%%%%%%%%%%%%%%%%%%%%%%%%%%%%%%%%%%%%%%%%%%%%%%%%%%%%%%%%%%%%%%
\begin{frame}[fragile]\frametitle{For Shapes like Sheet Metal \ldots}
\begin{center}
\includegraphics[width=\linewidth,keepaspectratio]{midcurve4}
\end{center}
Half the computation time, but similar accuracy
\end{frame}

%%%%%%%%%%%%%%%%%%%%%%%%%%%%%%%%%%%%%%%%%%%%%%%%%%%%%%%%%%%%%%%%%%%%%%%%%%%%%%%%%%
\begin{frame}[fragile]\frametitle{Midsurface is?}
\begin{center}
\includegraphics[width=0.9\linewidth,keepaspectratio]{midcurve5}
\end{center}
	\begin{itemize}
	\item Widely used for CAE of Thin-Walled parts
	\item Computation is challenging and still unsolved
	\end{itemize}
	
\end{frame}

%%%%%%%%%%%%%%%%%%%%%%%%%%%%%%%%%%%%%%%%%%%%%%%%%%%%%%%%%%%%%%%%%%%%%%%%%%%%%%%%%%
\begin{frame}[fragile]\frametitle{Getting Midsurface}

	\begin{itemize}
	\item Going on for decades \ldots
	\item Manually by offsetting and stitching, initially
	\item Many CAD-CAE packages give automatic option, but \ldots
	\end{itemize}
\begin{center}
\includegraphics[width=0.8\linewidth,keepaspectratio]{midcurve6}
\end{center}	
\end{frame}

%%%%%%%%%%%%%%%%%%%%%%%%%%%%%%%%%%%%%%%%%%%%%%%%%%%%%%%%%%%%%%%%%%%%%%%%%%%%%%%%%%
\begin{frame}[fragile]\frametitle{Look at the output}
\begin{center}
\includegraphics[width=\linewidth,keepaspectratio]{midcurve7}
\end{center}	
\end{frame}

%%%%%%%%%%%%%%%%%%%%%%%%%%%%%%%%%%%%%%%%%%%%%%%%%%%%%%%%%%%%%%%%%%%%%%%%%%%%%%%%%%
\begin{frame}[fragile]\frametitle{Can't tolerate gaps}

	\begin{itemize}
	\item We have thickness sampling, 
	\item To recreate-represent the original shape
	\item Input and output difference not desirable
	\end{itemize}
\end{frame}

%%%%%%%%%%%%%%%%%%%%%%%%%%%%%%%%%%%%%%%%%%%%%%%%%%%%%%%%%%%%%%%%%%%%%%%%%%%%%%%%%%
\begin{frame}[fragile]\frametitle{For a simple model like}
\begin{center}
\includegraphics[width=0.9\linewidth,keepaspectratio]{midcurve8}
\end{center}	
\end{frame}

%%%%%%%%%%%%%%%%%%%%%%%%%%%%%%%%%%%%%%%%%%%%%%%%%%%%%%%%%%%%%%%%%%%%%%%%%%%%%%%%%%
\begin{frame}[fragile]\frametitle{You get}
\begin{center}
\includegraphics[width=0.9\linewidth,keepaspectratio]{midcurve9}
\end{center}	
\end{frame}

%%%%%%%%%%%%%%%%%%%%%%%%%%%%%%%%%%%%%%%%%%%%%%%%%%%%%%%%%%%%%%%%%%%%%%%%%%%%%%%%%%
\begin{frame}[fragile]\frametitle{For a far simpler shape}
\begin{center}
\includegraphics[width=0.9\linewidth,keepaspectratio]{midcurve10}
\end{center}	
\end{frame}

%%%%%%%%%%%%%%%%%%%%%%%%%%%%%%%%%%%%%%%%%%%%%%%%%%%%%%%%%%%%%%%%%%%%%%%%%%%%%%%%%%
\begin{frame}[fragile]\frametitle{Current Quality}
\begin{center}
\includegraphics[width=0.9\linewidth,keepaspectratio]{midcurve11}
\end{center}
	\begin{itemize}
	\item Errors take weeks to correct for complex parts.
	\item But still preferred, due to vast savings time
	\item From Days to hours \ldots
	\end{itemize}
	
\end{frame}

%%%%%%%%%%%%%%%%%%%%%%%%%%%%%%%%%%%%%%%%%%%%%%%%%%%%%%%%%%%%%%%%%%%%%%%%%%%%%%%%%%
\begin{frame}[fragile]\frametitle{Midsurface Computation}

	\begin{itemize}
	\item Midsurface of a Patch is Midcurve of its profile extruded.
	\item So, it boils down to computing 1D midcurve of a 2D profile
	\end{itemize}
	
\begin{center}
\includegraphics[width=0.9\linewidth,keepaspectratio]{midcurve12}
\end{center}	
\end{frame}

%%%%%%%%%%%%%%%%%%%%%%%%%%%%%%%%%%%%%%%%%%%%%%%%%%%%%%%%%%%%%%%%%%%%%%%%%%%%%%%%%%
\begin{frame}[fragile]\frametitle{What is a Midcurve?}

	\begin{itemize}
	\item Midsurface : From 3D thin Solid to 2D Surface
	\item Midcurve : From 2D Profile to 1D Curve
	\end{itemize}
	
\begin{center}
\includegraphics[width=0.6\linewidth,keepaspectratio]{midcurve13}
\end{center}	
\end{frame}

%%%%%%%%%%%%%%%%%%%%%%%%%%%%%%%%%%%%%%%%%%%%%%%%%%%%%%%%%%%%%%%%%%%%%%%%%%%%%%%%%%
\begin{frame}[fragile]\frametitle{Many Approaches}

	\begin{itemize}
	\item More than 6 decades of research…
	\item Most CAD-CAE packages…
	\item Rule-based!! Heuristic!! Case-by-case basis!!
	\end{itemize}
	
	
\begin{center}
\includegraphics[width=0.9\linewidth,keepaspectratio]{midcurve14}
\end{center}	
\end{frame}

%%%%%%%%%%%%%%%%%%%%%%%%%%%%%%%%%%%%%%%%%%%%%%%%%%%%%%%%%%%%%%%%%%%%%%%%%%%%%%%%%%
\begin{frame}[fragile]\frametitle{When-What?}
\begin{center}
\includegraphics[width=\linewidth,keepaspectratio]{midcurve15}
\end{center}	
\end{frame}

%%%%%%%%%%%%%%%%%%%%%%%%%%%%%%%%%%%%%%%%%%%%%%%%%%%%%%%%%%%%%%%%%%%%%%%%%%%%%%%%%%
\begin{frame}[fragile]\frametitle{2017: My PhD Work: Rule-based}
\begin{center}
\includegraphics[width=\linewidth,keepaspectratio]{midcurve16}
\end{center}	
\end{frame}

%%%%%%%%%%%%%%%%%%%%%%%%%%%%%%%%%%%%%%%%%%%%%%%%%%%%%%%%%%%%%%%%%%%%%%%%%%%%%%%%%%
\begin{frame}[fragile]\frametitle{Limitations}

	\begin{itemize}
	\item Fully rule-based
	\item Need to adjust for new shapes
	\item So, not scalable
	\end{itemize}
\begin{center}

\includegraphics[width=0.3\linewidth,keepaspectratio]{midcurve17}
\end{center}
\end{frame}

% ── merged from midcurve_intro_ann.tex ─────────────────────────────────────────

%%%%%%%%%%%%%%%%%%%%%%%%%%%%%%%%%%%%%%%%%%%%%%%%%%%%%%%%%%%%%%%%%%%%%%%%%%%%%%%%%%
\begin{frame}[fragile]\frametitle{Idea}
\begin{center}

\includegraphics[width=0.3\linewidth,keepaspectratio]{midcurve18}

Can Neural Networks ``learn'' the dimension reduction transformation?
\end{center}
\end{frame}

%%%%%%%%%%%%%%%%%%%%%%%%%%%%%%%%%%%%%%%%%%%%%%%%%%%%%%%%%%%%%%%%%%%%%%%%%%%%%%%%%%
\begin{frame}[fragile]\frametitle{How?}

	\begin{itemize}
	\item Supply lots of training data of profiles and their corresponding midcurves and train.
	\item Then given an unseen profile, can Neural Network compute a midcurve, mimicking the original profile shape?
	\end{itemize}
\begin{center}

\includegraphics[width=0.3\linewidth,keepaspectratio]{midcurve18}
\end{center}
\end{frame}

%%%%%%%%%%%%%%%%%%%%%%%%%%%%%%%%%%%%%%%%%%%%%%%%%%%%%%%%%%%%%%%%%%%%%%%%%%%%%%%%%%
\begin{frame}[fragile]\frametitle{Midcurve by Neural network}
\begin{center}
\includegraphics[width=\linewidth,keepaspectratio]{midcurve19}
\end{center}
\end{frame}

%%%%%%%%%%%%%%%%%%%%%%%%%%%%%%%%%%%%%%%%%%%%%%%%%%%%%%%%%%%%%%%%%%%%%%%%%%%%%%%%%%
\begin{frame}[fragile]\frametitle{Midcurve : The Problem}

	\begin{itemize}
	\item {\bf Goal}: Given a 2D closed shape (closed polygon) find its midcurve (polyline, closed or open)
	\item {\bf Input}: set of points or set of connected lines, non-intersecting, simple, convex, closed polygon
	\item {\bf Output}: another set of points or set of connected lines, open/branched polygons possible
	\end{itemize}

\end{frame}

%%%%%%%%%%%%%%%%%%%%%%%%%%%%%%%%%%%%%%%%%%%%%%%%%%%%%%%%%%%%%%%%%%%%%%%%%%%%%%%%%%
\begin{frame}[fragile]\frametitle{Midcurve : Graph 2 Graph (The Ideal)}

  \begin{columns}[t]
    \begin{column}[T]{0.55\linewidth}
      \begin{itemize}
        \item {\bf Ideal framing}: Input profile graph $\rightarrow$ Output midcurve graph
        \item Both graphs have \textbf{different topologies}: different node/edge counts, positions
        \item Standard node-prediction and link-prediction GNNs assume similar input/output topology; not suitable here
        \item Direct graph-to-graph translation requires a coarsening / dimension-reduction formulation
        \item \alert{Phase III} of this work tackles this directly using a Graph Transformer
      \end{itemize}
    \end{column}
    \begin{column}[T]{0.45\linewidth}
      \centering
      \includegraphics[width=\linewidth,keepaspectratio]{midcurve_as_graph}
    \end{column}
  \end{columns}

\end{frame}

%%%%%%%%%%%%%%%%%%%%%%%%%%%%%%%%%%%%%%%%%%%%%%%%%%%%%%%%%%%%%%%%%%%%%%%%%%%%%%%%%%
\begin{frame}[fragile]\frametitle{Midcurve $==$ Dimension Reduction / Translation}

	\begin{itemize}
	\item Like PCA, find the principal curve that \textit{represents} the profile shape
	\item Left side (input): 2D Sketch Profile; Right side (output): 1D Midcurve
	\item \textbf{Not} an Auto-Encoder: input and output are \textit{different} shapes
	\item Variable-length input/output makes padding unsafe in geometry (any value is valid)
	\item Analogous to \textbf{graph summarization}: reduce a large graph to a smaller graph preserving structural essence, like text summarization keeps the meaning
	\end{itemize}
\begin{center}
\includegraphics[width=0.55\linewidth,keepaspectratio]{midcurve20}
\end{center}
\end{frame}

%%%%%%%%%%%%%%%%%%%%%%%%%%%%%%%%%%%%%%%%%%%%%%%%%%%%%%%%%%%%%%%%%%%%%%%%%%%%%%%%%%
\begin{frame}[fragile]\frametitle{Variable Size Encoder Decoder}

	\begin{itemize}
	\item Batches need fixed lengths
	\item Made fixed size by Padding.
	\end{itemize}
\begin{center}
\includegraphics[width=\linewidth,keepaspectratio]{midcurve21}
\end{center}
\end{frame}

%%%%%%%%%%%%%%%%%%%%%%%%%%%%%%%%%%%%%%%%%%%%%%%%%%%%%%%%%%%%%%%%%%%%%%%%%%%%%%%%%%
\begin{frame}[fragile]\frametitle{Why Padding Fails in Geometry}

	\begin{itemize}
	\item OK for NLP, say Machine Translations, where padding values like ``-1'' can be added along with other words (vectors or indices)
	\item But in Geometry, its not OK.
	\item Because any value can represent a Valid Input, even though we don't want it to be the input.
	\item Not all shapes can be serialised as sequences: Y-shapes and concentric rings (\textit{O}) cannot be traversed without lifting the pen, ruling out Seq2Seq models entirely
	\end{itemize}

\end{frame}

%%%%%%%%%%%%%%%%%%%%%%%%%%%%%%%%%%%%%%%%%%%%%%%%%%%%%%%%%%%%%%%%%%%%%%%%%%%%%%%%%%
\begin{frame}[fragile]\frametitle{Three Phase Strategy}

This trade-off motivates the \textbf{three-phase} strategy:
		\begin{itemize}
		\item \textbf{Phase I}: Raster Image $\rightarrow$ Image (encoder-decoder variants)
		\item \textbf{Phase II}: BRep JSON via LLMs, exact coords, current best performer
		\item \textbf{Phase III}: Geometry graph $\rightarrow$ graph, exact, native branches
		\end{itemize}

\end{frame}

%%%%%%%%%%%%%%%%%%%%%%%%%%%%%%%%%%%%%%%%%%%%%%%%%%%%%%%%%%%%%%%%%%%%%%%%%%%%%%%%%%
\begin{frame}[fragile]\frametitle{A Twist to the problem}

  \begin{columns}[t]
    \begin{column}[T]{0.4\linewidth}
      \centering
      \includegraphics[width=0.9\linewidth,keepaspectratio]{midcurve22}
    \end{column}
    \begin{column}[T]{0.6\linewidth}
	\begin{itemize}
	\item Till we get good variable size encoder decoder network for geometry…
	\item Decided to convert this Sequence 2 Sequence problem as Image 2 Image problem.
	\end{itemize}
    \end{column}
  \end{columns}
  \end{frame}
  
%%%%%%%%%%%%%%%%%%%%%%%%%%%%%%%%%%%%%%%%%%%%%%%%%%%%%%%%%%%%%%%%%%%%%%%%%%%%%%%%%%
\begin{frame}[fragile]\frametitle{A Twist to the problem}

	\begin{itemize}
	\item Input: Black \& White Image of 2D profile
	\item Output: Black \& White Image of 1D midcurve
	\end{itemize}
\begin{center}
\includegraphics[width=0.9\linewidth,keepaspectratio]{midcurve23}
\end{center}
\end{frame}

%%%%%%%%%%%%%%%%%%%%%%%%%%%%%%%%%%%%%%%%%%%%%%%%%%%%%%%%%%%%%%%%%%%%%%%%%%%%%%%%%%
\begin{frame}[fragile]\frametitle{Solves \ldots}
Problems of Geometric sequences

	\begin{itemize}
	\item Variable input/output sizes
	\item Loops need to be crossed
	\item Branches
	\end{itemize}
\begin{center}

\includegraphics[width=0.9\linewidth,keepaspectratio]{midcurve24}
\end{center}
\end{frame}
  
%%%%%%%%%%%%%%%%%%%%%%%%%%%%%%%%%%%%%%%%%%%%%%%%%%%%%%%%%%%%%%%%%%%%%%%%%%%%%%%%%%
\begin{frame}[fragile]\frametitle{Phase I: Raster Approximation}

	\begin{itemize}
	\item Images address both the \textbf{representation} and \textbf{variable-size} issues
	\item \alert{Big dilution}: True geometric shapes are \textit{vector} data; images here are \textit{raster}, so approximation has crept in
	\item Predicted output still needs post-processing to recover exact geometric (vector) form
	\end{itemize}

\end{frame}



%%%%%%%%%%%%%%%%%%%%%%%%%%%%%%%%%%%%%%%%%%%%%%%%%%%%%%%%%%%%%%%%%%%%%%%%%%%%%%%%%%
\begin{frame}[fragile]\frametitle{For Deep Learning}
	\begin{itemize}
	\item Need lots of data
	\item Had just few input output image pairs
	\item How to augment/populate large variations \ldots
	\end{itemize}
\end{frame}

% %%%%%%%%%%%%%%%%%%%%%%%%%%%%%%%%%%%%%%%%%%%%%%%%%%%%%%%%%%%%%%%%%%%%%%%%%%%%%%%%%%
% \begin{frame}[fragile]\frametitle{Data: Profile and Midcurve Pairs}
% \begin{center}
% \includegraphics[width=0.9\linewidth,keepaspectratio]{midcurve27}
% \end{center}
% \end{frame}

% %%%%%%%%%%%%%%%%%%%%%%%%%%%%%%%%%%%%%%%%%%%%%%%%%%%%%%%%%%%%%%%%%%%%%%%%%%%%%%%%%%
% \begin{frame}[fragile]\frametitle{Data: Few Pairs to Start With}
% \begin{center}
% \includegraphics[width=0.9\linewidth,keepaspectratio]{midcurve28}
% \end{center}
	% \begin{itemize}
	% \item For each shape, one input-output pair --- that's it
	% \item Must artificially multiply via geometric transformations
	% \end{itemize}
% \end{frame}

%%%%%%%%%%%%%%%%%%%%%%%%%%%%%%%%%%%%%%%%%%%%%%%%%%%%%%%%%%%%%%%%%%%%%%%%%%%%%%%%%%
\begin{frame}[fragile]\frametitle{Rasterization: Geometry to Image}
	\begin{itemize}
	\item Raw data is in vector format (list of $(x,y)$ points)
	\item Converted to fixed-size $100\times100$ bitmap via DrawSVG rasterization
	\end{itemize}
\begin{center}
\includegraphics[width=0.9\linewidth,keepaspectratio]{midcurve29}
\end{center}
\end{frame}

%%%%%%%%%%%%%%%%%%%%%%%%%%%%%%%%%%%%%%%%%%%%%%%%%%%%%%%%%%%%%%%%%%%%%%%%%%%%%%%%%%
\begin{frame}[fragile]\frametitle{Augmentation: Geometric Variations}

  \begin{columns}[t]
    \begin{column}[T]{0.4\linewidth}
      \centering
      \includegraphics[width=0.9\linewidth,keepaspectratio]{midcurve30}
    \end{column}
    \begin{column}[T]{0.6\linewidth}
	\begin{itemize}
	\item Base shapes: I, L, Plus, T
	\item Operations applied to both profile and midcurve:
	\begin{itemize}
		\item Translated
		\item Rotated
		\item Mirrored
		\item Mirrored + Translated
		\item Mirrored + Rotated
	\end{itemize}
	\item Total: 896 image pairs
	\end{itemize}
    \end{column}
  \end{columns}

\end{frame}

%%%%%%%%%%%%%%%%%%%%%%%%%%%%%%%%%%%%%%%%%%%%%%%%%%%%%%%%%%%%%%%%%%%%%%%%%%%%%%%%%%
\begin{frame}[fragile]\frametitle{Training Data Samples}
\begin{center}
\includegraphics[width=\linewidth,keepaspectratio]{midcurve31}
\end{center}
\end{frame}


%%%%%%%%%%%%%%%%%%%%%%%%%%%%%%%%%%%%%%%%%%%%%%%%%%%%%%%%%%%%%%%%%%%%%%%%%%%%%%%%%%
\begin{frame}[fragile]\frametitle{Reuse Image Encoder Decoder}
\begin{center}
\includegraphics[width=0.9\linewidth,keepaspectratio]{midcurve25}
\end{center}
\end{frame}

%%%%%%%%%%%%%%%%%%%%%%%%%%%%%%%%%%%%%%%%%%%%%%%%%%%%%%%%%%%%%%%%%%%%%%%%%%%%%%%%%%
\begin{frame}[fragile]\frametitle{For Dimension Reduction}
\begin{center}
\includegraphics[width=0.9\linewidth,keepaspectratio]{midcurve26}
\end{center}
\end{frame}

% ── Phase II: LLM / Text-based approach (overview) ────────────────────────────

%%%%%%%%%%%%%%%%%%%%%%%%%%%%%%%%%%%%%%%%%%%%%%%%%%%%%%%%%%%%%%%%%%%%%%%%%%%%%%%%%%
\begin{frame}[fragile]\frametitle{Phase II: Text as a Medium for Geometry}

  \begin{itemize}
    \item Phase I approximates geometry as pixel images: \alert{raster dilution}
    \item Phase II takes a different route: serialize geometry as \textbf{structured JSON text}
    \item LLMs can read and write JSON naturally: frame midcurve as \textbf{machine translation}:
      \begin{itemize}
        \item {\small Input: Profile geometry $\rightarrow$ BRep JSON text}
        \item {\small Output: Midcurve geometry $\leftarrow$ predicted BRep JSON text}
      \end{itemize}
    \item \textbf{Advantage over Phase I}: output is \textit{exact point coordinates}, not noisy pixels: no raster post-processing needed
    \item \textbf{Trade-off}: LLMs are probabilistic; output may occasionally need syntactic repair
    \item Models explored: Qwen2.5-7B, Gemma, Mistral (QLoRA fine-tuning); Nemotron-Mini-4B; few-shot prompting
  \end{itemize}

\end{frame}

%%%%%%%%%%%%%%%%%%%%%%%%%%%%%%%%%%%%%%%%%%%%%%%%%%%%%%%%%%%%%%%%%%%%%%%%%%%%%%%%%%
\begin{frame}[fragile]\frametitle{BRep: Boundary Representation in JSON}

  \textbf{Geometry encoded as a 3-level hierarchy} (2D adaptation of solid-modelling B-rep):

  \begin{columns}[t]
    \begin{column}[T]{0.48\linewidth}
      \begin{itemize}
        \item \textbf{Points}: ordered list of $(x,y)$ coordinates
        \item \textbf{Lines}: pairs of point indices: encodes \textit{connectivity}
        \item \textbf{Segments}: chains of line indices: encodes \textit{continuous open curves}; \textbf{multiple Segments} represent junctions in T- and Plus-shapes
      \end{itemize}
      \vspace{0.5em}
      {\small Key invariant: only \textit{coordinates} change under geometric transforms; topology (Lines, Segments) stays fixed}
    \end{column}
    \begin{column}[T]{0.52\linewidth}
      {\footnotesize\textbf{I-shape Profile:}}\\[0.15em]
      {\footnotesize
        \texttt{Points:} $(5,5),(10,5),(10,20),(5,20)$\\
        \texttt{Lines:}  $[0,1],[1,2],[2,3],[3,0]$\\
        \texttt{Segs:}   $[0,1,2,3]$
      }\\[0.6em]
      {\footnotesize\textbf{I-shape Midcurve:}}\\[0.15em]
      {\footnotesize
        \texttt{Points:} $(7.5,5),(7.5,20)$\\
        \texttt{Lines:}  $[0,1]$\\
        \texttt{Segs:}   $[0]$
      }\\[0.6em]
      {\footnotesize\textbf{T-shape Midcurve (branched):}}\\[0.15em]
      {\footnotesize
        \texttt{Segs:} $[0],[1],[2]$ \quad (3 arms = 3 Segments)\\
        \textit{Sequences cannot encode this; BRep can.}
      }
    \end{column}
  \end{columns}

\end{frame}

%%%%%%%%%%%%%%%%%%%%%%%%%%%%%%%%%%%%%%%%%%%%%%%%%%%%%%%%%%%%%%%%%%%%%%%%%%%%%%%%%%
\begin{frame}[fragile]\frametitle{Phase II: How LLMs Predict Midcurves}

  \begin{columns}[t]
    \begin{column}[T]{0.55\linewidth}
      \begin{itemize}
        \item \textbf{Data}: 993 BRep examples: 4 base shapes (I, L, T, Plus) $\times$ geometric augmentations (rotate, scale, translate, mirror)
        \item \textbf{Framing}: text-to-text machine translation
          \begin{itemize}
            \item {\small Input: \texttt{Profile\_brep} JSON string}
            \item {\small Output: \texttt{Midcurve\_brep} JSON string}
          \end{itemize}
        \item \textbf{Sub-approaches}:
          \begin{itemize}
            \item {\small Few-shot prompting (ChatGPT, Bard, Claude)}
            \item {\small QLoRA fine-tuning (4-bit quantised Low-Rank Adaptation, trains ${<}1\%$ of params): Qwen2.5-7B, Gemma, Mistral}
            \item {\small Nemotron-Mini-4B: HF SFTTrainer / Unsloth / few-shot}
          \end{itemize}
      \end{itemize}
    \end{column}
    \begin{column}[T]{0.45\linewidth}
      \centering
      \includegraphics[width=\linewidth,keepaspectratio]{chatgpt_lines_15dec2023_prompt}
    \end{column}
  \end{columns}

\end{frame}

%%%%%%%%%%%%%%%%%%%%%%%%%%%%%%%%%%%%%%%%%%%%%%%%%%%%%%%%%%%%%%%%%%%%%%%%%%%%%%%%%%
\begin{frame}[fragile]\frametitle{Phase II: Results}

  \begin{itemize}
    \item \textbf{Qwen2.5-7B (QLoRA fine-tuning)}: \alert{current best performer}:
      \begin{itemize}
        \item MAE = 0.78, RMSE = 1.24, Hausdorff distance = 2.1
        \item \textbf{PSR (Profile Shape Reconstruction) = 98\%} on 100-sample test set
      \end{itemize}
    \item \textbf{Nemotron-Mini-4B (2-shot, no fine-tuning)}:
      \begin{itemize}
        \item PSR = 85.7\%, topology accuracy = 0.83
        \item Fits in $\sim$2\,GB VRAM: deployable on consumer hardware
      \end{itemize}
    \item \textbf{Key advantage}: output is exact \textit{point coordinates}: no pixel clean-up
    \item \textbf{Limitation}: branched shapes (T, Plus) require well-formed multi-Segment JSON output
  \end{itemize}

\end{frame}

% ── Phase III: Pure geometry / graph-based approach (overview) ─────────────────

%%%%%%%%%%%%%%%%%%%%%%%%%%%%%%%%%%%%%%%%%%%%%%%%%%%%%%%%%%%%%%%%%%%%%%%%%%%%%%%%%%
\begin{frame}[fragile]\frametitle{Phase III: Back to Pure Geometry: Graph to Graph}

  \begin{columns}[t]
    \begin{column}[T]{0.55\linewidth}
      \begin{itemize}
        \item Recall the ideal: input profile graph $\rightarrow$ output midcurve graph
        \item Phase III implements this directly: \textbf{no pixels, no text tokens}
        \item \textbf{Input}: polygon as graph: nodes = corner points with $(x,y)$; edges = adjacent corners
        \item \textbf{Output}: midcurve as graph: different node/edge counts and positions
        \item \textbf{Advantage}: exact coordinates, native branch handling, zero post-processing
        \item Branches (T, Plus) are naturally captured in the graph topology
      \end{itemize}
    \end{column}
    \begin{column}[T]{0.45\linewidth}
      \centering
      \includegraphics[width=\linewidth,keepaspectratio]{midcurve_as_graph}
    \end{column}
  \end{columns}

\end{frame}

%%%%%%%%%%%%%%%%%%%%%%%%%%%%%%%%%%%%%%%%%%%%%%%%%%%%%%%%%%%%%%%%%%%%%%%%%%%%%%%%%%
\begin{frame}[fragile]\frametitle{Phase III: Graph Transformer Architecture}

  \begin{itemize}
    \item \textbf{Core challenge}: graph-to-graph translation with \textit{different} input/output topology: not standard node/link prediction
    \item \textbf{Non-auto-regressive Graph Transformer} (trained from scratch):
      \begin{itemize}
        \item LaplacianPE positional encoding $\rightarrow$ TransformerConv layers $\rightarrow$ TopKPooling
        \item TopKPooling acts as graph \textit{coarsening}: the geometric dimension-reduction step
        \item Loss: Chamfer distance (point-cloud matching) + BCE (adjacency matrix)
      \end{itemize}
    \item \textbf{Second variant}: pretrained Graphormer (HuggingFace) fine-tuned in two phases (frozen head $\rightarrow$ joint fine-tuning)
    \item Operates on raw $(x,y)$ coordinates: no rasterization, no JSON encoding
    \item Status: \alert{implemented, evaluation ongoing}
  \end{itemize}

\end{frame}

% ── Three-phase comparison + Conclusion ───────────────────────────────────────

%%%%%%%%%%%%%%%%%%%%%%%%%%%%%%%%%%%%%%%%%%%%%%%%%%%%%%%%%%%%%%%%%%%%%%%%%%%%%%%%%%
\begin{frame}[fragile]\frametitle{Three Phases: A Progression Toward the Ideal}

\begin{center}
\small
\begin{tabular}{@{}p{2.0cm}p{2.4cm}p{2.4cm}p{2.4cm}@{}}
\toprule
 & \textbf{Phase I} & \textbf{Phase II} & \textbf{Phase III} \\
 & Image-based      & LLM / Text        & Geometry           \\
\midrule
Medium       & Bitmap (raster)     & BRep JSON text      & Graph $(x,y)$      \\
Exact coords & No (pixels)         & Yes                 & Yes                \\
Branches     & No                  & Yes (Segments)      & Yes (native)       \\
Best model   & CNN encoder-decoder & Qwen2.5-7B QLoRA    & Graph Transformer  \\
Key metric   & Visual quality      & PSR = 98\%          & Ongoing eval.      \\
Status       & Fully impl.         & \alert{Best performer}    & Implemented        \\
\bottomrule
\end{tabular}
\end{center}

\end{frame}

%%%%%%%%%%%%%%%%%%%%%%%%%%%%%%%%%%%%%%%%%%%%%%%%%%%%%%%%%%%%%%%%%%%%%%%%%%%%%%%%%%
\begin{frame}[fragile]\frametitle{Conclusion}

  \begin{itemize}
    \item Midcurve is a \textbf{geometric dimension-reduction} problem: 2D profile $\rightarrow$ 1D skeleton (graph $\rightarrow$ smaller graph)
    \item \textbf{Phase I} (image-based): encoder-decoder on raster images --- feasible; CNN encoder-decoder is the best primary variant; rasterization trade-off limits precision
    \item \textbf{Phase II} (LLM / text-based): BRep JSON + LLM fine-tuning: \alert{current best performer}: PSR = 98\%, exact coordinates, branching via Segments
    \item \textbf{Phase III} (geometry-based): Graph Transformer on polygon graphs: the ideal formulation; no approximation, native branch handling; evaluation ongoing
    \item \textbf{Next frontier}: generalise to arbitrary 2D thin polygons; extend the graph-to-graph paradigm toward 3D midsurface computation
  \end{itemize}

\begin{center}
\alert{\large Graph-to-Graph is the ultimate goal for MidcurveNN.}
\end{center}

\end{frame}
